problem:  
Let \((M^m,g_M)\) be a complete, noncompact Riemannian manifold with nonnegative Ricci curvature, and let \((N^n,g_N)\) be a complete, simply connected Riemannian manifold with nonpositive sectional curvature satisfying  
\[
K_N(x) \le -\,f(r(x)), \qquad r(x)=d_N(p,x),
\]  
where \(p\in N\) is fixed and \(f:[0,\infty)\to[0,\infty)\) is smooth, nondecreasing, and \(f(r)\to0\) as \(r\to\infty\).  

Assume \(u:M\to N\) is a harmonic map of finite total energy:  
\[
E(u)=\frac12\int_M |\nabla u|^2 <\infty.
\]  

Known results imply constancy of such maps when \(K_N\le -c<0\).  
Now suppose instead that \(f(r)\approx C(1+r)^{-\alpha}\) for some \(\alpha>0\).  

**Question:**  
For which exponents \(\alpha\) (possibly depending on \(m=\dim M\) and the volume growth of \(M\)) must every finite-energy harmonic map \(u:M\to N\) be constant?  
Equivalently, determine the *sharp critical exponent* \(\alpha_c(m)\) such that  
- if \(f(r)\gg (1+r)^{-\alpha_c}\), rigidity holds (every map is constant),  
- if \(f(r)\ll (1+r)^{-\alpha_c}\), nontrivial finite-energy harmonic maps may exist.  

Why is it a "good" problem:  
This problem captures the sought-after boundary between rigidity and flexibility for harmonic maps, governed by the *rate at which negative curvature decays at infinity*. It translates geometric curvature decay into quantitative analytic rigidity, merging insights from geometric analysis, potential theory, and the theory of elliptic PDEs. Its apparent simplicity—asking for a critical exponent in a power-law decay—conceals formidable conceptual depth: the answer likely depends on subtle interactions between the curvature of \(N\), the volume growth of \(M\), and the integrability of the energy density. Thus it bridges global Riemannian geometry and harmonic map theory through a precise parametric question. Determining (or even conjecturing) the value of such a critical decay exponent could unify several Liouville-type results under one asymptotic principle, providing a sharp, quantitative advance in our understanding of vanishing theorems for harmonic maps.